% ============================================================================
\chapter{Hardverska platforma i FreeRTOS port}
\label{ch:platforma}
% ============================================================================

U ovom poglavlju opisuje se osnova na kojoj po\v{c}iva cijeli rad. Najprije se
predstavlja razvojna plo\v{c}a i meki procesor na kojem se izvr\v{s}ava program,
zatim prilago\dj{}eni port operativnog sistema FreeRTOS za taj procesor, te
na\v{c}in na koji su organizovani memorija i ulazno-izlazni registri sistema na
\v{c}ipu. Na kraju poglavlja prikazuje se osnovni dokaz da operativni sistem
zaista radi na hardveru, jer bez te potvrde nijedan kasniji rezultat ne bi imao
smisla. Ovo poglavlje predstavlja temelj za razumijevanje kasnijih poglavlja u
kojima se aplikacija razdvaja na zadatke i uklapa u stroge resurse plo\v{c}e.

\section{Tang Nano 9K i meki procesor PicoRV32}
\label{sec:platforma-hardver}

Razvojna plo\v{c}a Tang Nano 9K zasnovana je na programabilnom logi\v{c}kom polju
proizvo\dj{}a\v{c}a Gowin, modela GW1NR-9C. Plo\v{c}a se taktuje vanjskim
oscilatorom od 27 megaherca i raspola\v{z}e ograni\v{c}enim skupom logi\v{c}kih
elemenata i blokovske memorije, \v{s}to je kasnije klju\v{c}no za uklapanje
sistema. Na plo\v{c}i je putem sinteze digitalnog sklopa realizovan meki
procesor PicoRV32, jednostavna jezgra zasnovana na otvorenoj skupovnoj
arhitekturi RISC-V \citep{picorv32}. Procesor je konfigurisan za cjelobrojni
skup instrukcija sa mno\v{z}enjem, \v{s}to se u alatima ozna\v{c}ava kao
\foreignlanguage{english}{\emph{rv32im}}.

PicoRV32 je vi\v{s}eciklusna jezgra, \v{s}to zna\v{c}i da svaka instrukcija
tro\v{s}i vi\v{s}e takta. Takva izvedba je mala i \v{s}tedljiva po pitanju
logike, ali sporija po izvr\v{s}enoj instrukciji u odnosu na cjevovodne procesore.
Ta osobina je va\v{z}na za kasnije tuma\v{c}enje izmjerenih vrijednosti, jer se
broj takta po zamjeni konteksta ne mo\v{z}e direktno tuma\v{c}iti kao broj
izvr\v{s}enih instrukcija.

\section{Prilago\dj{}eni port operativnog sistema FreeRTOS}
\label{sec:platforma-port}

FreeRTOS je operativni sistem za rad u realnom vremenu otvorenog koda, dostupan
za veliki broj arhitektura \citep{freertos}. Za potrebe ovog sistema kori\v{s}ten
je prilago\dj{}eni port jezgre verzije V10.5.1 za PicoRV32. Raspore\dj{}iva\v{c}
periodi\v{c}no dobija takt iz vanjskog broja\v{c}a koji generi\v{s}e prekid na
liniji broj pet, podesen na frekvenciju od sto herca. Promjena zadatka, odnosno
zamjena konteksta, pokre\'ce se programskim prekidom na liniji broj \v{s}est.

Va\v{z}na posebnost ovog porta je obavezna upotreba prevodila\v{c}ke opcije koja
isklju\v{c}uje takozvanu relaksaciju, \foreignlanguage{english}{\emph{-mno-relax}}.
Razlog je taj \v{s}to PicoRV32 u kori\v{s}tenoj konfiguraciji pri ulasku u prekid
mo\v{z}e prepisati registre globalnog i nitnog pokaziva\v{c}a, pa se program ne
smije oslanjati na globalni pokaziva\v{c}. Ova naizgled sitna postavka klju\v{c}na
je za ispravan rad sistema i predstavlja primjer pa\v{z}nje koju zahtijeva rad
na niskom nivou.

\section{Memorijska i adresna organizacija sistema na \v{c}ipu}
\label{sec:platforma-memorija}

Sistem na \v{c}ipu razdvaja prostor programa i prostor podataka. Program se
nalazi u zasebnoj memoriji, takozvanom programskom prozoru, dok se podaci i stog
nalaze u radnoj memoriji veli\v{c}ine \v{s}esnaest kilobajta. Pristup
perifernim jedinicama ostvaruje se preslikavanjem u memoriju
(\foreignlanguage{english}{\emph{Memory-Mapped Input/Output}}, MMIO), gdje svaka
periferija zauzima svoj raspon adresa. Serijska jedinica, brojac, op\v{s}ti
ulazi i izlazi te video jedinica imaju odvojene adresne opsege. Pojednostavljena
slika ove organizacije prikazana je na slici \ref{fig:memorija}.

\begin{figure}[H]
\centering
\begin{tikzpicture}[
    blok/.style={draw, rounded corners, minimum width=5.6cm, minimum height=0.85cm, align=center, font=\small},
    node distance=0.35cm
]
\node[blok] (uart) {Serijska jedinica (UART) -- 0x02000000};
\node[blok, below=of uart] (timer) {Brojac / takt (IRQ5) -- 0x03000000};
\node[blok, below=of timer] (gpio) {Op\v{s}ti ulazi/izlazi (dugmad, LED) -- 0x04000000};
\node[blok, below=of gpio] (video) {Video jedinica (znakovi, mapa) -- 0x05000000};
\node[blok, below=of video, fill=black!5] (ram) {Radna memorija (podaci, stog) -- 0x00000000, 16 KB};
\node[blok, below=of ram, fill=black!5] (rom) {Programski prozor (kod) -- 0x00100000};
\end{tikzpicture}
\caption{Pojednostavljena adresna i memorijska organizacija sistema na \v{c}ipu.}
\label{fig:memorija}
\end{figure}

Razdvajanje programa i podataka, kao i veli\v{c}ine pojedinih memorija, kasnije
postaje sredi\v{s}nji problem pri uklapanju cijele igre sa video izlazom, jer
svaka od tih memorija tro\v{s}i ograni\v{c}enu blokovsku memoriju plo\v{c}e.

\section{Dokaz pokretanja na hardveru}
\label{sec:platforma-boot}

Prije bilo kakvog portanja aplikacije, potvrdjeno je da operativni sistem uop\v{s}te
radi na plo\v{c}i. Sintetizovan je sklop sa serijskim izlazom, a u njega je
ugra\dj{}en jednostavan program sa tri zadatka. Nakon upisa sklopa u memoriju
plo\v{c}e, preko serijske veze na brzini od sto petnaest hiljada bita u sekundi
dobijeni su o\v{c}ekivani ispisi zadataka i njihovo smjenjivanje. Time je
pokazano da takt raspore\dj{}iva\v{c}a, programski prekid za zamjenu konteksta i
me\dj{}usobno isklju\v{c}ivanje pristupa serijskoj jedinici rade ispravno.

\section{Zaklju\v{c}ak poglavlja}
\label{sec:platforma-zakljucak}

U poglavlju je predstavljena hardverska osnova rada. Opisani su plo\v{c}a Tang
Nano 9K i meki procesor PicoRV32, prilago\dj{}eni port operativnog sistema
FreeRTOS sa taktom iz vanjskog brojaca i zamjenom konteksta putem programskog
prekida, te organizacija memorije i perifernih jedinica. Potvr\dj{}eno je da
sistem radi na stvarnom hardveru. Ova osnova omogu\'cava da se u narednom
poglavlju pristupi razdvajanju aplikacije na zasebne zadatke.
